\chapter{Gaussian tools for arbitrary profiles}\label{sec:gauss}

The results of \cref{sec:fnd} concern the law of \cref{subsec:fnd-law},
and the sharper ones hold only at a minimizing profile. In
\cref{sec:sc,sec:c7,sec:c7-endpoint}, we also need laws whose profile
minimizes no functional, and laws whose auxiliary field has different
variances in different coordinates. In this section, we introduce a
class of generalized laws that contains these cases
(\cref{def:gauss-law}), and we prove two Gaussian tools for them.

The first tool is \cref{lem:fnd-posterior}, which shows that the
martingale identities and the covariance order hold for every generalized
law. It also bounds the variance of averages over
levels of the squared conditional mean on a set of coordinates. We note
that this bound holds for every profile and that its proof does not use
contact identities. We use it in \cref{sec:sc} for the
scalar law, which is the generalized law with $N=1$, field variance $d$,
no couplings, and root mass $s=0$. We also use it in
\cref{sec:c7,sec:c7-endpoint} for the cavity laws, which have diffusion
matrix $\operatorname{diag}(cI_n,b(u))$ and root mass $s=\lambda/\beta$.

The second tool is \cref{lem:fnd-path-heat}, a version of the Gaussian
heat identity for normalized expectations. Here we take expectations of
functionals of finitely many Gaussian
variables and continuous Gaussian paths, and we hold the profile measure
fixed. We use the lemma in
\cref{sec:sc} to differentiate the scalar law under transports of its
levels, and in \cref{sec:c7,sec:c7-endpoint} to differentiate the cavity
laws in the cavity parameter. In this section, we use only the definitions of
\cref{sec:fnd}, the arguments of \cite[Section~2]{AL}, and the proof of
\cref{prop:fnd-contact}, which we adapt.

Note that some letters have a local meaning in this section. The letter
$\omega$ denotes Gaussian couplings, while in \cref{sec:c7}, $\omega$
denotes the
weights of a tree. In \cref{sec:gauss-heat}, the derivative measures $M$
and $Q$ are unrelated to $M=Nc$, to the quantile function
\eqref{eq:fnd-quantile}, and to the vector $M_q$ of \cref{sec:c7}.

\section{Generalized laws and a posterior variance bound}
\label{sec:gauss-posterior}

In \cref{sec:c7}, we decouple one spin from the others. The remaining
spins then evolve under a law whose profile need not minimize any
functional and whose auxiliary field has unequal variances. The
following class of laws contains both situations.

\begin{definition}[Generalized law]\label{def:gauss-law}
Let $\Lambda=\operatorname{diag}(\vartheta_1,\ldots,\vartheta_N)$ with
$\vartheta_i>0$, and let $(\kappa_{ij})_{1\le i\le j\le N}$ be real
numbers; $\Lambda$ and $\kappa$ are deterministic. Let
$\omega=(\omega_{ij})_{1\le i\le j\le N}$ be independent standard
Gaussian variables, and write $\E_{\rm ref}$ for expectation under their
law. Let $\mu=s\delta_0+\nu$ be either an admissible profile in the
sense of \cref{sec:fnd-functional}, so that $0<s\le1/2$, or, with
$s=0$, an arbitrary deterministic probability measure $\mu=\nu$ on
$[0,1]$. Let $a(q)=\mu([0,q])$, and let $U$ solve
\begin{equation}\label{eq:fnd-general-pde}
 \partial_qU=-\frac12\bigl(\Tr(\Lambda\nabla^2U)
      +a(q)\nabla U^T\Lambda\nabla U\bigr),\qquad
 U(1,x;\omega)=\log\sum_{\sigma\in\{-1,1\}^N}
     e^{\mathcal H_\sigma(\omega,x)},
\end{equation}
where
$\mathcal H_\sigma(\omega,x)=\sum_{i\le j}\kappa_{ij}\omega_{ij}\sigma_i\sigma_j+x\cdot\sigma$.
The \emph{generalized law} of $(\Lambda,\kappa,\mu)$ is the joint law of
$(\omega,X,\sigma)$ built in three stages.
\begin{enumerate}[label=(\roman*)]
\item The Gaussian law of $\omega$ is tilted by the density
$e^{sU(0,0;\omega)}/\E_{\rm ref}e^{sU(0,0;\omega)}$.
\item Given $\omega$, the field solves
$\dd X_q=a(q)\Lambda\nabla U(q,X_q;\omega)\dd q+\Lambda^{1/2}\dd B_q$
with $X_0=0$, where $B$ is a standard Brownian motion in $\R^N$.
\item Given $\omega$ and $X$, the spin $\sigma$ is drawn with
probabilities proportional to $e^{\mathcal H_\sigma(\omega,X_1)}$.
\end{enumerate}
As in \cref{subsec:fnd-law}, we write $\Pp$ and $\E$ for this law, and
we define $\calG_q=\sigma(\omega,X_r:0\le r\le q)$,
$m_q=\E[\sigma\mid\calG_q]$, $C_q=\Cov(\sigma\mid\calG_q)$, and
$H_q=\nabla_x^2U(q,X_q;\omega)$.
\end{definition}

When $s=0$, the outer density in (i) is identically one. Recall that a
Hopf--Cole step of mass zero is an
ordinary Gaussian expectation. The
fractional-moment functionals of \cref{sec:fnd} still have $s>0$. The
equation \eqref{eq:fnd-general-pde} has a solution for every such
profile, which is obtained from step profiles by the Hopf--Cole recursion and
uniform approximation, as we show in Step~6 of the proof of
\cref{lem:fnd-posterior}. With $\Lambda=cI$, with $\omega_{ij}=g_{ij}$
and $\kappa_{ij}=\beta\sqrt{t/N}$ for $i<j$, and with $\omega_{ii}=z_i$
and $\kappa_{ii}=\beta\sqrt{t/(2N)}$, we have
$\mathcal H_\sigma(\omega,x)=\beta\sqrt t\,\sigma^TW\sigma/2+x\cdot\sigma$.
In this case, the generalized law is the law of \cref{subsec:fnd-law}. The scalar
law of \cref{sec:sc} is the generalized law with $N=1$, $\Lambda=d$,
$\kappa=0$, and $s=0$. The cavity laws of \cref{sec:c7} are generalized
laws with $\Lambda=\operatorname{diag}(cI_n,b(u))$ and
independent Gaussian couplings with deterministic variances.

In the next lemma, we adapt the argument of the proof of
\cite[Lemma~2.3]{OpenAIFluctuationScale}, which is also used in the proof
of \cite[Lemma~2.1]{OpenAILimitingLaw}. That argument combines joint
convexity in the Gaussian drivers and a spin source, the inverse-Hessian
Brascamp--Lieb inequality, and a gauge symmetry that removes the
conditioning on the terminal spin. We use
it for averages of squared conditional means on selected coordinates, and
for diagonal anisotropic laws.

\begin{lemma}[Averaged self-overlap under an arbitrary profile]
\label{lem:fnd-posterior}
Let $\mu=s\delta_0+\nu$ be an admissible profile, and consider either
its joint law from \cref{subsec:fnd-law}, for some $0<\rho\le1$, or a
generalized law; a generalized law also allows $s=0$ and an arbitrary
probability measure $\mu=\nu$. Then
$m_q=\nabla U(q,X_q)=\E[\sigma\mid\calG_q]$, $m_0=0$, and
$a(q)C_q\preceq H_q\preceq C_q$ for every $q\in[0,1]$. Fix a set of
$n\ge1$ coordinates, let $m_q^-$ be the restriction of $m_q$ to these
coordinates, and let $H_q^{--}$ and $C_q^{--}$ be the corresponding
principal submatrices of $H_q$ and $C_q$. For every deterministic
$k\in L^2(\nu)$,
\begin{equation}\label{eq:fnd-posterior}
 \begin{aligned}
 \Var\left(\int k(q)\frac{|m_q^-|^2}{n}\,\nu(\dd q)\right)
 &\le\frac4{n^2}\int k(q)^2
       \E[(m_q^-)^TH_q^{--}m_q^-]\,\nu(\dd q)\\
 &\le\frac4{n^2}\int k(q)^2
       \E[(m_q^-)^TC_q^{--}m_q^-]\,\nu(\dd q).
 \end{aligned}
\end{equation}
\end{lemma}

Recall that the law of \cref{subsec:fnd-law} is a generalized law, so
that it suffices to prove the lemma for generalized laws. The proof also
gives the
martingale differential
\begin{equation}\label{eq:gauss-martingale}
 \dd m_q=H_q\Lambda^{1/2}\dd B_q\qquad(0\le q\le1)
\end{equation}
for every generalized law, as we will see in Steps~1 and~6 below.

\begin{proof}[Proof of \cref{lem:fnd-posterior}]
The proof has six steps. We first prove the martingale identities and
the covariance order for step profiles. Next, still for step profiles,
we write the law in terms of independent Gaussian coordinates. We
compare the curvature of the posterior density with the gradient of the
squared conditional mean, and we apply the Brascamp--Lieb inequality
with the terminal spin fixed. We then remove the conditioning on the
spin by a gauge symmetry. Finally, we pass to general profiles and
general $k$.

\paragraph{Step 1: martingale and covariance order.}
Let $\mu$ be a step profile. We adapt Step~1 of the proof of
\cref{prop:fnd-contact}, replacing $cI$ by $\Lambda$. The generator of
the field is
$\Tr(\Lambda\nabla^2)/2+a(q)(\Lambda\nabla U)\cdot\nabla$.
Differentiating \eqref{eq:fnd-general-pde} in $x$ and applying It\^o's
formula along the field on each interval between grid points, we find
that
\[
 \dd\,\nabla U(q,X_q;\omega)=H_q\Lambda^{1/2}\dd B_q,\qquad
 \dd H_q=-a(q)H_q\Lambda H_q\dd q+\nabla H_q\Lambda^{1/2}\dd B_q ,
\]
where the last term is the contraction of the third spatial derivative
of $U$ with $\Lambda^{1/2}\dd B_q$. The stochastic integrals are martingales,
since the spatial derivatives are bounded at fixed $N$. The
gradient $\nabla U(1,X_1;\omega)$ is the Gibbs mean of the terminal spin.
By the tower property, we have $\nabla U(q,X_q;\omega)=\E[\sigma\mid\calG_q]=m_q$,
which proves \eqref{eq:gauss-martingale} for step profiles. Using
$C_1=H_1$, the martingale isometry, and the equation for $H$, we find
\[
 \begin{gathered}
 C_q=\E[H_1\mid\calG_q]+\int_q^1\E[H_u\Lambda H_u\mid\calG_q]\dd u,
 \\
 H_q=\E[H_1\mid\calG_q]+\int_q^1a(u)\E[H_u\Lambda H_u\mid\calG_q]\dd u .
 \end{gathered}
\]
Note that the matrices $H_u\Lambda H_u$ are positive semidefinite and
that $a(q)\le a(u)\le1$ for $u\ge q$. Subtracting the two identities as in
Step~1 of the proof of \cref{prop:fnd-contact}, we obtain
$a(q)C_q\preceq H_q\preceq C_q$. Note that
$\mathcal H_{-\sigma}(\omega,x)=\mathcal H_\sigma(\omega,-x)$. Replacing
$\sigma$ by $-\sigma$ in the terminal sum, we see that
$U(1,\cdot\,;\omega)$ is even. The equation preserves this symmetry,
which implies that $U(q,\cdot\,;\omega)$ is even and that $m_0=0$.

\paragraph{Step 2: Gaussian coordinates and the posterior density.}
Let $0=q_0<q_1<\cdots<q_K=1$ and $a=a_\ell$ on $[q_\ell,q_{\ell+1})$,
with $a_\ell\ge s\ge0$. Then the cumulative function is constant between
grid points,
and the free measure $\nu$ is carried by the grid points, with
$\nu(\{q_0\})=a_0-s$, $\nu(\{q_\ell\})=a_\ell-a_{\ell-1}$ for
$1\le\ell\le K-1$, and $\nu(\{q_K\})=1-a_{K-1}$. Let
$z_1,\ldots,z_K$ be independent standard Gaussian vectors in $\R^N$, let
\[
 X_{q_\ell}=\sum_{h=1}^\ell(q_h-q_{h-1})^{1/2}\Lambda^{1/2}z_h
 \qquad(0\le\ell\le K),
\]
and let $Z$ collect the coordinates of $\omega$ and of
$z_1,\ldots,z_K$. We write $U_q(Z)=U(q,X_q;\omega)$ at the grid points.
We recall that \eqref{eq:fnd-general-pde} is solved by the Hopf--Cole
recursion
$U_{q_\ell}=a_\ell^{-1}\log\E_{z_{\ell+1}}e^{a_\ell U_{q_{\ell+1}}}$ when
$a_\ell>0$, and by $U_{q_\ell}=\E_{z_{\ell+1}}U_{q_{\ell+1}}$ when
$a_\ell=0$. We follow the construction of \cite[(2.1)--(2.3)]{AL}. We claim
that, restricted to $\omega$, to the field at the grid points, and to the
spin, the generalized law is the law of
$(\omega,(X_{q_\ell})_\ell,\sigma)$ when $(Z,\sigma)$ has the density
\[
 \frac{\exp\{\mathcal H_\sigma(\omega,X_1)
                    -\int U_q(Z)\,\nu(\dd q)\}}
      {\E_{\rm ref}e^{sU(0,0;\omega)}}
\]
with respect to the standard Gaussian law of $Z$ and counting measure on
$\{-1,1\}^N$. Indeed, the outer tilt contributes
$sU_{q_0}$ to the exponent, the transition from $q_\ell$ to
$q_{\ell+1}$ contributes $a_\ell(U_{q_{\ell+1}}-U_{q_\ell})$, and the
terminal Gibbs probability contributes $\mathcal H_\sigma(\omega,X_1)-U_{q_K}$.
In the sum, we see that the coefficients of
$U_{q_0},U_{q_1},\ldots,U_{q_K}$ are
\[
 s-a_0,\quad a_0-a_1,\quad\ldots,\quad a_{K-2}-a_{K-1},\quad a_{K-1}-1,
\]
which are the negatives of the masses of $\nu$ at the grid points.
A transition of mass zero has likelihood one, and the calculation also
holds at $s=0$, where the denominator equals one. For a fixed spin
$\sigma$, we find that the negative logarithm of the conditional density
of $Z$ is, up to a constant, the posterior potential
\[
 \Phi_\sigma(Z)=\frac12|Z|^2-\mathcal H_\sigma(\omega,X_1)
           +\int U_q(Z)\,\nu(\dd q).
\]
Note that $\mathcal H_\sigma(\omega,X_1)$ is linear in $Z$ and that each
$U_q$ is convex in $Z$. To see the convexity, observe that the terminal value
in \eqref{eq:fnd-general-pde} is a log-sum-exp of linear functions of
$(\omega,x)$, that each Hopf--Cole step of positive mass preserves
convexity by H\"older's inequality, and that Gaussian expectation at
mass zero preserves it as well. Since $X_q$ is linear in $Z$,
we obtain
\[
 \nabla_Z^2\Phi_\sigma=\mathcal Q=I+\int\mathcal A_q\,\nu(\dd q)\succeq I,
 \qquad
 \mathcal A_q=\nabla_Z^2U_q\succeq0 .
\]
The matrix $\mathcal Q$ does not depend on $\sigma$.

\paragraph{Step 3: joint source curvature.}
Let $\iota\colon\R^n\to\R^N$ embed a vector in the selected coordinates,
and let $\mathcal U_q(Z,y)=U(q,X_q(Z)+\iota y;\omega)$ for
$y\in\R^n$. This function is jointly convex in $(Z,y)$, since $U(q,x;\omega)$ is
convex in $(\omega,x)$ by Step~2 and the map
$(Z,y)\mapsto(\omega,X_q(Z)+\iota y)$ is linear. We write
$\mathcal B_q=\nabla_Z\nabla_y\mathcal U_q(Z,0)$, a matrix with $n$
columns. Since $\nabla_y\mathcal U_q(Z,0)=m_q^-$ and
$\nabla_y^2\mathcal U_q(Z,0)=H_q^{--}$, joint convexity and the chain
rule give
\begin{equation}\label{eq:fnd-source-curvature}
 \begin{pmatrix}
  \mathcal A_q&\mathcal B_q\\
  \mathcal B_q^T&H_q^{--}
 \end{pmatrix}\succeq0,\qquad
 \nabla_Z|m_q^-|^2=2\mathcal B_qm_q^- .
\end{equation}
For a vector $w$ and $r\in\R$, we evaluate this positive semidefinite
form at $(rw,m_q^-)$. The resulting quadratic polynomial in $r$ is
nonnegative. This implies that its discriminant is nonpositive:
\[
 |w^T\mathcal B_qm_q^-|^2
 \le(w^T\mathcal A_qw)\,(m_q^-)^TH_q^{--}m_q^- .
\]
The same statements hold at $q=0$, where $m_0^-=0$.

\paragraph{Step 4: the Brascamp--Lieb inequality.}
Let $Y=n^{-1}\int k(q)|m_q^-|^2\,\nu(\dd q)$, a finite sum over grid
points, and let $b=\int k(q)\mathcal B_qm_q^-\,\nu(\dd q)$, so that
$\nabla_ZY=2b/n$ by \eqref{eq:fnd-source-curvature}. For every vector $w$,
using the last display, the Cauchy--Schwarz
inequality for the integral against $\nu$, and
$\int\mathcal A_q\,\nu(\dd q)=\mathcal Q-I\preceq\mathcal Q$, we obtain
\[
 \begin{aligned}
 |w^Tb|^2&\le\Bigl(\int w^T\mathcal A_qw\,\nu(\dd q)\Bigr)
             \int k(q)^2(m_q^-)^TH_q^{--}m_q^-\,\nu(\dd q)\\
 &\le(w^T\mathcal Qw)\int k(q)^2(m_q^-)^TH_q^{--}m_q^-\,\nu(\dd q).
 \end{aligned}
\]
Taking $w=\mathcal Q^{-1}b$ and dividing by $b^T\mathcal Q^{-1}b$ when it is
positive, we find that
\[
 b^T\mathcal Q^{-1}b\le\int k(q)^2(m_q^-)^TH_q^{--}m_q^-\,\nu(\dd q);
\]
if $b^T\mathcal Q^{-1}b=0$, this inequality holds trivially. Fix $\sigma$. The
potential $\Phi_\sigma$ is smooth, its Hessian $\mathcal Q$ is at least
$I$, and it has Gaussian tails. By the Brascamp--Lieb inequality
\cite{BrascampLieb1976}, in the form
\cite[Proposition~2.2, (2.14)]{AL}, we have
$\Var(Y\mid\sigma)\le\E[\nabla_ZY^T\mathcal Q^{-1}\nabla_ZY\mid\sigma]$,
and we conclude that
\[
 \Var(Y\mid\sigma)\le\frac4{n^2}
   \E\Bigl[\int k(q)^2(m_q^-)^TH_q^{--}m_q^-\,\nu(\dd q)\Bigm|\sigma\Bigr].
\]

\paragraph{Step 5: removing the conditioning.}
For $\tau\in\{-1,1\}^N$, we consider the map
\[
 \sigma_i\mapsto\tau_i\sigma_i,\qquad
 \omega_{ij}\mapsto\tau_i\tau_j\omega_{ij},\qquad
 z_{h,i}\mapsto\tau_iz_{h,i}.
\]
This map preserves the Gaussian reference law. Because $\Lambda$ is diagonal,
it multiplies each $X_q$ coordinatewise by $\tau$. Further,
$\mathcal H_{\tau\sigma}$ evaluated at the transformed $(\omega,x)$
equals $\mathcal H_\sigma(\omega,x)$, and the map leaves every $U_q$
invariant. It therefore preserves the joint density of Step~2. It maps
$m_{q,i}$ to $\tau_im_{q,i}$ and $(H_q)_{ij}$ to $\tau_i\tau_j(H_q)_{ij}$,
so that $Y$ and $(m_q^-)^TH_q^{--}m_q^-$ are invariant. The map acts
transitively on spin configurations, which implies that the conditional
law of these quantities given $\sigma$ does not depend on $\sigma$. In
particular, $\E[Y\mid\sigma]$ is constant, and the law of total variance
implies that $\Var(Y)=\E\Var(Y\mid\sigma)$. Combining this with Step~4, we
obtain the first inequality in
\eqref{eq:fnd-posterior} for step profiles. The second inequality follows
from $H_q^{--}\preceq C_q^{--}$,
since principal submatrices preserve the order $H_q\preceq C_q$.

\paragraph{Step 6: general profiles and general $k$.}
We approximate $\mu$ by step profiles with the same root mass, as in
Step~4 of the proof of \cref{prop:fnd-contact} and
\cite[Section~2.6]{AL}. That is, we push $\nu$ forward to finite meshes that
contain $0$, $1$, and any prescribed finite set of levels, with mesh size
tending to zero. The argument of \cite[Section~2.6]{AL} uses only that
the spins are bounded, that spatial derivatives of fixed order are
bounded uniformly over step profiles at fixed $N$, that the solution is
Lipschitz in the profile in $L^1$, that the outer tilt has Gaussian
domination, and that the disorder, the Brownian motion, and the terminal
spin can be coupled. These facts carry over to a bounded diagonal
$\Lambda$ and to couplings $\kappa_{ij}\omega_{ij}$. For the Lipschitz
bound, let $U$ and $\widetilde U$ be the solutions for the cumulative
functions $a$ and $\tilde a$, and let $\varphi=U-\widetilde U$. Subtracting
the two equations, we find
\[
 \partial_q\varphi+\frac12\Tr(\Lambda\nabla^2\varphi)
 +\frac a2(\nabla U+\nabla\widetilde U)^T\Lambda\nabla\varphi
 =-\frac{a-\tilde a}{2}\,\nabla\widetilde U^T\Lambda\nabla\widetilde U,
 \qquad \varphi(1,\cdot\,;\omega)=0 .
\]
Since each coordinate of $\nabla\widetilde U$ is a conditional spin
mean, we have $\nabla\widetilde U^T\Lambda\nabla\widetilde U\le\Tr\Lambda$. By
the backward representation of $\varphi$, we have
\[
 \sup_{x,\omega}|\varphi(q,x;\omega)|\le\frac{\Tr\Lambda}2\int_q^1|a(u)-\tilde a(u)|\dd u
 \le\frac N2\max_i\vartheta_i\,\|a-\tilde a\|_{L^1([0,1])}.
\]
This replaces the constant $cN/2$ of \cite[(2.19)]{AL}. It follows that
the ratio of the outer densities for $\mu$ and for an approximating step
profile converges to one uniformly in $\omega$. At $s=0$, the outer
density is identically one, and the approximation requires no division
by a root mass.

Under the common coupling of \cite[Section~2.6.1]{AL}, the processes
$X$, $m$, $H$, and $C$ for the step profiles converge to those for $\mu$,
uniformly in $q$, in probability. Then the identities of Step~1, including
\eqref{eq:gauss-martingale},
pass to the limit, where for \eqref{eq:gauss-martingale} we use the It\^o
isometry and the bounded
convergence of $H$. For continuous $k$, the integrands
$q\mapsto k(q)|m_q^-|^2$ and $q\mapsto k(q)^2(m_q^-)^TH_q^{--}m_q^-$ are
bounded and continuous along each path. From any subsequence we may
extract a further subsequence on which the process convergence is
almost surely uniform in $q$. On that subsequence the integrands
converge uniformly, and their integrals then converge by weak
convergence of the deterministic free measures. The integrands are bounded, the ratio of the outer densities converges
to one uniformly in $\omega$, and, since $U(0,0;\omega)$ grows at most
linearly in $\omega$, the outer densities are dominated, uniformly along
the approximation, by a function that is integrable under the Gaussian
reference law. This allows passage to the expectations and variances. Thus \eqref{eq:fnd-posterior} holds for continuous $k$.
Next, we
show that both sides of \eqref{eq:fnd-posterior} are continuous
in $k\in L^2(\nu)$. Since $|m_q^-|^2\le n$, we have
\[
 \Bigl\|\frac1n\int\bigl(k_1(q)-k_2(q)\bigr)|m_q^-|^2\,\nu(\dd q)\Bigr\|_{L^2(\Pp)}
 \le\nu([0,1])^{1/2}\|k_1-k_2\|_{L^2(\nu)},
\]
which shows that the left side of \eqref{eq:fnd-posterior} is
continuous in $k$. Further,
$(m_q^-)^TH_q^{--}m_q^-\le(m_q^-)^TC_q^{--}m_q^-\le n\|C_q\|_{\rm op}\le nN$,
which implies that the right sides change by at most
$4n^{-1}N\|k_1-k_2\|_{L^2(\nu)}\|k_1+k_2\|_{L^2(\nu)}$. Since continuous
functions are dense in $L^2(\nu)$, this proves the lemma.
\end{proof}

\section{Differentiating normalized Gaussian path expectations}
\label{sec:gauss-heat}

In \cref{sec:sc,sec:c7,sec:c7-endpoint}, we differentiate normalized
expectations under Gaussian field laws while keeping their profile
measures fixed. The parameter changes the covariance of the Gaussian
variables. In
\cref{sec:sc}, it changes the clock of the field path, and in
\cref{sec:c7,sec:c7-endpoint}, it changes the variance of a row of
couplings and the clock of one field coordinate. The densities that we
encounter are exponentials of
integrals of smooth functions of the field against the profile measure,
which may have atoms. We use the following version of the Gaussian heat
identity, for finitely many Gaussian variables and continuous field
paths. Its finite-dimensional input is the classical formula for the
derivative of a Gaussian expectation with respect to the covariance. A
general version, under hypotheses
different from ours, is proved in \cite[Theorem~1]{Voigtlaender2021Price}.
In the proof, we reduce to
finite Gaussian vectors by Bernstein approximation of the paths, while we
keep the profile measure, including its atoms, exact.

\paragraph{Setting.}
Let $I\subset\R$ be a nondegenerate compact interval with interior $I^\circ$, and let
$z\in I$ be the parameter. Under a reference probability measure, with
expectation $\E_{\rm ref}$, we consider the following random
variables.
\begin{enumerate}[label=(\roman*)]
\item \emph{Finite coordinates}
$G^{\rm fin}_z=(v_k(z)^{1/2}\xi_k)_{1\le k\le n_0}$, where
$\xi_1,\ldots,\xi_{n_0}$ are independent standard Gaussian variables and
$v_k\colon I\to[0,\infty)$.
\item \emph{Field paths} on a fixed finite rooted tree. Each edge $e$
carries the level interval $[0,1]$, after an affine change of variable,
and $d$ components. For each edge $e$ and component $j\le d$, let
$B_{e,j}$ be a standard Brownian motion, and let
$T_{e,j,z}\colon[0,1]\to[0,\infty)$ be a deterministic nondecreasing clock
such that $T_{e,j,z}(0)=0$. The Brownian motions for different edges
and components are mutually independent. The path on the edge $e$ is
\[
 G_z(e,j,v)=G_z(e',j,1)+B_{e,j}\bigl(T_{e,j,z}(v)\bigr),
 \qquad 0\le v\le1,
\]
where $e'$ is the parent edge of $e$, and the first term is $0$ if $e$
starts at the root.
\item A finite family $\omega$ of independent standard Gaussian variables that do not
depend on $z$ and are independent of the variables in (i) and (ii), and
a variable $\sigma$ that ranges over a nonempty finite set and is integrated
against counting measure.
\end{enumerate}
The variables $\xi_k$ and the Brownian motions $B_{e,j}$ are mutually
independent. Thus the finite coordinates and the field paths are
independent under the reference law. This setting includes all the
scalar and cavity applications below.
We write $\mathcal I$ for the index space formed by the finite
coordinates $1\le k\le n_0$ and the field points and components
$(e,j,v)$. It is a compact metric space, the disjoint union of a finite
set and finitely many copies of $[0,1]$. The Gaussian coordinates
$G_z=(G_z^{\rm fin},G_z^{\rm path})$ form a random element of the space
$\mathcal C$ of functions on $\mathcal I$ that are continuous on each
copy of $[0,1]$, with the norm $|G^{\rm fin}|+\|G^{\rm path}\|_\infty$.
We denote their covariance by $K_z(i,i')=\E_{\rm ref}[G_z(i)G_z(i')]$ for
$i,i'\in\mathcal I$. For example, for one edge and one component,
$K_z(v,v')=T_z(\min\{v,v'\})$.

Let $\ell_z(G)=\ell_z(G,\omega,\sigma)$ and
$\Phi_z(G)=\Phi_z(G,\omega,\sigma)$ be real functions of $z\in I$,
$G\in\mathcal C$, $\omega$, and $\sigma$. We write $\dot\ell_z$ and
$\dot\Phi_z$ for their derivatives in $z$, with $(G,\omega,\sigma)$
fixed, and $\dot K_z$ for the derivative of $K_z$ in $z$. For a
continuous kernel $A$ on $\mathcal I^2$ and a finite signed measure $B$
on $\mathcal I^2$, we write $\langle A,B\rangle=\iint A(i,i')B(\dd i,\dd i')$,
and $S\otimes P$ denotes the product of two signed measures on
$\mathcal I$.

\begin{lemma}[Normalized Gaussian path heat identity]
\label{lem:fnd-path-heat}
In the setting above, assume the following.
\begin{enumerate}[label=\textup{(A\arabic*)}]
\item \emph{Variances and clocks.} Each $v_k$ and each $T_{e,j,z}(v)$ is
continuous in $z\in I$, jointly with $v$ for the clocks, and
continuously differentiable in $z\in I^\circ$. For $z\in I^\circ$, each
nonconstant $v_k$ is positive and each clock is strictly increasing in
$v$. Coordinates with $v_k\equiv0$ are omitted.
\item \emph{Covariance.} The covariance $K_z(i,i')$ is differentiable in
$z\in I^\circ$, and the maps $(z,i,i')\mapsto K_z(i,i')$ and
$(z,i,i')\mapsto\dot K_z(i,i')$ are jointly continuous on
$I\times\mathcal I^2$, where $\dot K_z$ at the endpoints of $I$ denotes
the continuous extension.
\item \emph{Functionals.} For each $\sigma$, the functions
$\ell_z(G,\omega,\sigma)$, $\Phi_z(G,\omega,\sigma)$,
$\dot\ell_z(G,\omega,\sigma)$, and $\dot\Phi_z(G,\omega,\sigma)$ are jointly
continuous in $(z,G,\omega)\in I\times\mathcal C\times\R^{n_1}$,
where $n_1$ is the number of auxiliary Gaussian coordinates, with
$\dot\ell_z$ and $\dot\Phi_z$ extended continuously to the endpoints of $I$.
The functions $\Phi_z$, $\dot\Phi_z$, and
$\dot\ell_z$ are bounded uniformly in all variables, and there exists
$C<\infty$ such that
\[
 |\ell_z(G)|\le C\bigl(1+|\omega|+|G^{\rm fin}|+\|G^{\rm path}\|_\infty\bigr).
\]
\item \emph{Derivative measures.} For fixed $(z,\omega,\sigma)$, the maps
$G\mapsto\ell_z(G)$ and $G\mapsto\Phi_z(G)$ are twice continuously
differentiable on $\mathcal C$, and their derivatives are represented by
finite signed measures
\[
 S=D\ell_z,\qquad M=D^2\ell_z,\qquad
 P=D\Phi_z,\qquad Q=D^2\Phi_z
\]
on $\mathcal I$ and $\mathcal I^2$, in the sense that, for example,
\[
 D\ell_z(G)[h]=\int h\dd S,\qquad
 D^2\ell_z(G)[h,h']=\iint h(i)h'(i')\,M(\dd i,\dd i')
 \qquad(h,h'\in\mathcal C).
\]
For fixed $(\omega,\sigma)$, the four measures depend continuously on
$(z,G)$ in the total-variation norm, and their total variations are
bounded uniformly in all variables.
\end{enumerate}
Let $Z_z=\E_{\rm ref}e^{\ell_z(G_z)}$, where $\E_{\rm ref}$ includes the
sum over $\sigma$, and let
$\E_z\Phi=Z_z^{-1}\E_{\rm ref}\bigl[e^{\ell_z(G_z)}\Phi\bigr]$. Then, for
every $z\in I$,
\begin{align}
 \frac\dd{\dd z}\E_z\Phi_z={}&
 \E_z\dot\Phi_z+\Cov_z(\Phi_z,\dot\ell_z)\nonumber\\
 &+\frac12\E_z\Bigl\langle\dot K_z,
 Q+S\otimes P+P\otimes S
 +(\Phi_z-\E_z\Phi_z)(M+S\otimes S)\Bigr\rangle,
 \label{eq:fnd-path-heat}
\end{align}
where the functionals and measures are evaluated at $G_z$. The
derivative is one-sided at the endpoints of $I$, including endpoints at
which a finite-coordinate variance vanishes. The same differentiation
gives
\begin{equation}\label{eq:fnd-path-normalizer}
 \frac\dd{\dd z}\log Z_z
 =\E_z\dot\ell_z+\frac12\E_z\langle\dot K_z,M+S\otimes S\rangle.
\end{equation}
\end{lemma}

\begin{remark}\label{rem:gauss-local}
We show in the proof that the right sides of \eqref{eq:fnd-path-heat} and
\eqref{eq:fnd-path-normalizer} are continuous in $z\in I$, so that
$z\mapsto\E_z\Phi_z$ and $z\mapsto\log Z_z$ are continuously
differentiable on $I$. The assumptions hold, in particular, for the
smooth local functionals that we use in
\cref{sec:sc,sec:c7,sec:c7-endpoint},
integrated against a fixed finite measure. For instance, let $\varpi$ be
a finite measure on the path part of $\mathcal I$, possibly with atoms,
and let $\ell_z(G)=-\int V_z(i,G(i))\,\varpi(\dd i)$, where $V_z(i,x)$
is twice continuously differentiable in $x$, and $\partial_xV_z$ and
$\partial_x^2V_z$ are bounded and jointly continuous in $(z,i,x)$. Then
\[
 S=-\partial_xV_z(\cdot,G(\cdot))\,\varpi,\qquad
 M=-\operatorname{diag}_{\#}\bigl(\partial_x^2V_z(\cdot,G(\cdot))\,\varpi\bigr),
\]
where $\operatorname{diag}_{\#}$ is the pushforward under
$i\mapsto(i,i)$. The total variation of the difference of the scores at
$G$ and $G'$ is at most
$\varpi(\mathcal I)\sup_i|\partial_xV_z(i,G(i))-\partial_xV_z(i,G'(i))|$,
and the same holds for the Hessians with $\partial_x^2V_z$. We note that
these functions are uniformly continuous on compact sets
and that $G$ is bounded, which gives the continuity in (A4) for this
term. An atom of $\varpi$ produces an atom of $M$ on
the diagonal, and its contribution to $\langle\dot K_z,M\rangle$
involves $\dot K_z$ on the diagonal.
\end{remark}

\begin{proof}[Proof of \cref{lem:fnd-path-heat}]
The proof has seven steps. We first prove the identities for finite
Gaussian vectors. We then approximate the paths by their Bernstein
polynomials and couple the approximations on one probability space.
Next, we show that the derivative measures and the covariance kernels
converge, and we establish a common integrable envelope and a lower
bound for the normalizers. Finally, we pass to the limit in the interior
of $I$ and treat the endpoints.

\paragraph{Step 1: the finite-dimensional identity.}
We begin with a finite Gaussian vector. Let $J$ be an open interval,
and let $Y_z$, $z\in J$, be a centered Gaussian vector in $\R^p$ whose
covariance $\Sigma_z$ is continuously differentiable in $z$. We allow the
matrix $\Sigma_z$ to be
singular. Let $F_z\colon\R^p\to\R$ be continuously
differentiable in $z$ and twice continuously differentiable in
$y\in\R^p$, and assume that $F_z$, $\partial_zF_z$, $\nabla F_z$, and
$\nabla^2F_z$ are bounded by $Ce^{C|y|}$, locally uniformly in $z$. We
claim that
\begin{equation}\label{eq:gauss-finite-heat}
 \frac\dd{\dd z}\E F_z(Y_z)=\E\,\partial_zF_z(Y_z)
 +\frac12\E\bigl\langle\dot\Sigma_z,\nabla^2F_z(Y_z)\bigr\rangle .
\end{equation}
Let $\varepsilon>0$, let $\zeta$ be a standard Gaussian vector in
$\R^p$, and let $A_z=(\Sigma_z+\varepsilon I)^{1/2}$. The matrix square
root is smooth on positive definite matrices, so $A_z$ is continuously
differentiable in $z$. By dominated convergence, we may differentiate
$\E F_z(A_z\zeta)$ under the expectation. The chain rule gives, besides
$\E\,\partial_zF_z(A_z\zeta)$, the term
$\E[\nabla F_z(A_z\zeta)\cdot\dot A_z\zeta]$. Using Gaussian integration
by parts, $\E[\zeta_lg(\zeta)]=\E[\partial_lg(\zeta)]$, we find that
this term equals $\E\Tr\bigl(\dot A_z\nabla^2F_z(A_z\zeta)A_z\bigr)$.
Since $\nabla^2F_z$ is symmetric and
$\dot A_zA_z+A_z\dot A_z=\dot\Sigma_z$, this is
$\E\langle\dot\Sigma_z,\nabla^2F_z(A_z\zeta)\rangle/2$. We integrate
this identity over $[z_1,z_2]\subset J$ and let $\varepsilon\downarrow0$.
The vector $A_z\zeta$ has the law of $Y_z+\sqrt\varepsilon\zeta'$, where
$\zeta'$ is a standard Gaussian vector independent of $Y_z$. Since the
growth bounds give domination, we obtain
the integrated form of
\eqref{eq:gauss-finite-heat}. Realizing $Y_z=\Sigma_z^{1/2}\zeta$, which
is continuous in $z$, we see by the same domination that the right side
of \eqref{eq:gauss-finite-heat} is continuous in $z$. The fundamental
theorem of calculus then gives
\eqref{eq:gauss-finite-heat}. This is the classical formula. Note that we
never invert $\Sigma_z$.

We will apply \eqref{eq:gauss-finite-heat} to $e^{\ell_z}$ and to
$e^{\ell_z}\Phi_z$, composed with a linear map from $\R^p$ to
$\mathcal C$, as in Step~2. By the product rule, we have
$D e^{\ell}=e^{\ell}S$, $D^2e^{\ell}=e^{\ell}(M+S\otimes S)$, and
\[
 D^2(e^{\ell}\Phi)=e^{\ell}\bigl[Q+S\otimes P+P\otimes S+\Phi(M+S\otimes S)\bigr].
\]
Inserting these derivatives into \eqref{eq:gauss-finite-heat} for the
numerator $\E_{\rm ref}[e^{\ell_z}\Phi_z]$ and the denominator
$\E_{\rm ref}e^{\ell_z}$, and using the quotient rule, we obtain
\eqref{eq:fnd-path-heat} for a finite Gaussian vector. Observe that the
centering $\Phi_z-\E_z\Phi_z$ comes from the
derivative of the denominator.
Dividing the derivative of the denominator by the denominator gives
\eqref{eq:fnd-path-normalizer}.

\paragraph{Step 2: Bernstein approximation and kernel convergence.}
We use the Bernstein approximation introduced in
\cite{Bernstein1912}. Let $m\ge1$ and $v_l=l/m$ for $0\le l\le m$. For a
continuous function $f$ on $[0,1]$, the Bernstein polynomial is
\[
 \Pi_mf(v)=\sum_{l=0}^mb_{m,l}(v)f(v_l),\qquad
 b_{m,l}(v)=\binom ml v^l(1-v)^{m-l}.
\]
The operator $\Pi_m$ is linear, and $\Pi_mf(0)=f(0)$ and
$\Pi_mf(1)=f(1)$. The weights are nonnegative and sum to one, which
gives $\|\Pi_mf\|_\infty\le\|f\|_\infty$. Further, for every $\delta>0$, we have
\[
 \|\Pi_mf-f\|_\infty\le\omega_f(\delta)+\frac{\|f\|_\infty}{2m\delta^2},
\]
where $\omega_f$ is the modulus of continuity of $f$. This follows from
Chebyshev's inequality for a binomial variable with $m$ trials, whose
variance is at most $m/4$. In particular, $\Pi_mf\to f$ uniformly
over any bounded equicontinuous family of functions $f$. The same holds for the
tensor product $\Pi_m\otimes\Pi_m$ acting on continuous functions on
$[0,1]^2$, with the bound $\omega_f(\delta)+\|f\|_\infty/(m\delta^2)$.

On each edge and component, we let $\Pi_m$ act on the path, and we
leave the finite coordinates unchanged. We denote the resulting operator
on $\mathcal C$ also by $\Pi_m$. Then $\Pi_mG_z$ is a linear function of
the finite Gaussian vector $Y_z$ formed by the finite coordinates and
the sampled values $G_z(e,j,v_l)$. A value at a vertex of the tree may
appear in several coordinates of $Y_z$. Step~1 allows such repeated
coordinates, since it does not require an invertible covariance. The
covariance of $Y_z$ consists of the values of $K_z$ at the sampled
points. By (A2), it is continuously differentiable on $I^\circ$, with
derivative given by the values of $\dot K_z$. The growth conditions of
Step~1 hold by (A3) and (A4), since $\|\Pi_my\|_{\mathcal C}\le2|y|$.

We set
\[
 \ell^{(m)}_z=\ell_z(\Pi_mG_z),\qquad
 \Phi^{(m)}_z=\Phi_z(\Pi_mG_z),\qquad
 Z^{(m)}_z=\E_{\rm ref}e^{\ell^{(m)}_z},
\]
and we normalize the approximate moment
$\E^{(m)}_z\Phi^{(m)}_z=(Z^{(m)}_z)^{-1}\E_{\rm ref}[e^{\ell^{(m)}_z}\Phi^{(m)}_z]$
by its own normalizer $Z^{(m)}_z$. By the chain rule, the derivative of
$\ell_z(\Pi_my)$ in the sampled coordinate $(e,j,v_l)$ is
$\int b_{m,l}\dd S$, where $b_{m,l}$ is placed on the edge $e$ and the
component $j$. We obtain the second derivatives in the same way from
$M$. We use Step~1 for fixed $(\omega,\sigma)$, and we then integrate
over $\omega$ and sum over $\sigma$. The domination of Step~5 justifies
this. We conclude that \eqref{eq:fnd-path-heat} and
\eqref{eq:fnd-path-normalizer} hold for the approximations at every
$z\in I^\circ$, with $\E_z$ replaced by $\E^{(m)}_z$ and with the
functionals and
measures evaluated at $\Pi_mG_z$. In these identities, $\dot K_z$ is
replaced by the kernel
\[
 \dot K_z^{(m)}(v,w)
  =\sum_{l,l'}b_{m,l}(v)b_{m,l'}(w)\dot K_z(v_l,v_{l'})
\]
on each pair of edges and components. Blocks between finite coordinates
are unchanged, and blocks between a finite coordinate $k$ and a path use
the same sampling in their path variable:
\[
 \dot K_z^{(m)}(k,(e,j,v))=\sum_lb_{m,l}(v)\dot K_z(k,(e,j,v_l)).
\] By (A2),
the kernels $K_z$ and $\dot K_z$ are uniformly continuous on the
compact set $I\times\mathcal I^2$. Hence the families
$\{K_z\}_{z\in I}$ and $\{\dot K_z\}_{z\in I}$ are bounded and
equicontinuous. Using the Bernstein bound, we obtain
\begin{equation}\label{eq:fnd-path-kernel-limit}
 \sup_{z\in I}\|\dot K_z^{(m)}-\dot K_z\|_\infty\longrightarrow0,
 \qquad
 \sup_{z\in I}\|K_z^{(m)}-K_z\|_\infty\longrightarrow0,
\end{equation}
where $K_z^{(m)}$, the covariance of $\Pi_mG_z$, is defined in the same
way.

\paragraph{Step 3: coupling.}
In Step~1, we used only the law of $Y_z$ at each fixed $z$. To pass to
the limit, we use the realization of the setting, with one Brownian motion
$B_{e,j}$ for each edge and component, evaluated at the varying clocks.
The finite coordinates are $v_k(z)^{1/2}\xi_k$ for fixed
$\xi_k$, independent of these Brownian motions as in the setting.
By (A1), the map $(z,v)\mapsto G_z(e,j,v)$ is almost surely
continuous on the compact set $I\times[0,1]$. It is therefore uniformly
continuous, and the family $\{G_z(e,j,\cdot)\}_{z\in I}$ is
equicontinuous. By Step~2, we have
$\sup_{z\in I}\|\Pi_mG_z-G_z\|_{\mathcal C}\to0$ almost surely. Let
$\bar T$ be the maximum of the clocks over $I\times[0,1]$, with
$\bar T=0$ if there are no field paths. Let
$\bar v=\max(\{0\}\cup\{v_k(z):1\le k\le n_0,\ z\in I\})$, so that
$\bar v=0$ if there are no finite coordinates. Since the value of a path
at a point of the edge $e$ is a sum of
increments of the Brownian motions on $e$ and its ancestor edges, we find
that
\[
 \sup_{z\in I}\|\Pi_mG_z\|_{\mathcal C}\le\sup_{z\in I}\|G_z\|_{\mathcal C}
 \le\bar v^{1/2}|\xi|+\sum_{e,j}\sup_{0\le\tau\le\bar T}|B_{e,j}(\tau)| .
\]

\paragraph{Step 4: convergence of the derivative measures.}
By (A4) and Step~3, we see that the measures $S$, $M$, $P$, and $Q$
evaluated at $\Pi_mG_z$ converge in total variation to the same measures evaluated
at $G_z$, almost surely, for every $z\in I$. Combining this with
\eqref{eq:fnd-path-kernel-limit} and the uniform bounds on the total
variations, we find that every contraction in the identity of Step~2
converges. For example, writing $M^{(m)}$ for $M$ evaluated at $\Pi_mG_z$,
we have
\[
 \bigl|\langle\dot K_z^{(m)},M^{(m)}\rangle-\langle\dot K_z,M\rangle\bigr|
 \le\|\dot K_z^{(m)}-\dot K_z\|_\infty\|M^{(m)}\|_{\rm TV}
 +\|\dot K_z\|_\infty\|M^{(m)}-M\|_{\rm TV},
\]
and for product measures we use
$\|S'\otimes P'-S\otimes P\|_{\rm TV}\le\|S'-S\|_{\rm TV}\|P'\|_{\rm TV}
+\|S\|_{\rm TV}\|P'-P\|_{\rm TV}$. This includes the diagonal
contributions of atoms, as in
\cref{rem:gauss-local}, since only the field values are approximated
and $\dot K^{(m)}_z(v,v)\to\dot K_z(v,v)$ uniformly.

\paragraph{Step 5: domination and the normalizers.}
By (A3) and Step~3, we have $e^{\ell^{(m)}_z}\le\mathcal W$ for all $m$
and $z$, where
\[
 \mathcal W=\exp\Bigl\{C\Bigl(1+|\omega|+\bar v^{1/2}|\xi|
 +\sum_{e,j}\sup_{0\le\tau\le\bar T}|B_{e,j}(\tau)|\Bigr)\Bigr\}
\]
and $C$ is the constant of (A3). All differentiated terms in the
identity of Step~2 are products of a weight with quantities bounded by
(A3), (A4), and \eqref{eq:fnd-path-kernel-limit}. Hence a constant
multiple of $\mathcal W$ is a common envelope. It is integrable to every fixed
positive power, because the suprema of Brownian motion over a bounded
interval have Gaussian tails, and the sum over $\sigma$ is finite. We now
show that the approximate normalizers have a common positive
lower bound. We retain one value $\sigma_0$ of the counting variable.
Using Jensen's inequality and the bound on $|\ell_z|$ in (A3), we obtain
\[
 Z^{(m)}_z\ge\E_{\rm ref}e^{\ell_z(\Pi_mG_z,\omega,\sigma_0)}
 \ge\exp\bigl(\E_{\rm ref}\,\ell_z(\Pi_mG_z,\omega,\sigma_0)\bigr)\ge e^{-C'},
\]
where the expectations are over the Gaussian variables only, and $C'$
does not depend on $m$ or $z$.

\paragraph{Step 6: passage to the limit in the interior.}
Let $\mathcal E_m(z)=\E^{(m)}_z\Phi^{(m)}_z$, and let $\Psi_m(z)$ be the
right side of the identity of Step~2. Let $[z_1,z_2]\subset I^\circ$.
By Step~2 and the fundamental theorem of calculus, we have
\[
 \mathcal E_m(z_2)-\mathcal E_m(z_1)=\int_{z_1}^{z_2}\Psi_m(z)\dd z .
\]
By Steps~3--5 and dominated convergence, we see that
$\mathcal E_m(z)\to\E_z\Phi_z$
and $\Psi_m(z)\to\Psi(z)$ for all $z\in I$, where
$\Psi(z)$ is the right side of \eqref{eq:fnd-path-heat}. Here we used the
lower bound on the normalizers and the continuity of $\ell_z$, $\Phi_z$,
$\dot\ell_z$, and $\dot\Phi_z$ in $G$. The
functions $\Psi_m$ are bounded uniformly in $m$ and $z$. Using dominated
convergence in $z$, we obtain
\[
 \E_{z_2}\Phi_{z_2}-\E_{z_1}\Phi_{z_1}=\int_{z_1}^{z_2}\Psi(z)\dd z .
\]
The same argument with $z_k\to z$ in place of $m\to\infty$ shows that
$\Psi$ and $z\mapsto\E_z\Phi_z$ are continuous on $I$. Indeed, by the
coupling, $G_{z_k}\to G_z$ uniformly, the measures converge in total
variation by (A4), $\dot K_{z_k}\to\dot K_z$ uniformly by (A2), and the
envelope of Step~5 dominates. By the fundamental theorem of calculus, we
conclude that
\eqref{eq:fnd-path-heat} holds on $I^\circ$. The same argument with
$\Phi\equiv1$, applied to $\log Z^{(m)}_z$ instead of a ratio, gives
\eqref{eq:fnd-path-normalizer} on $I^\circ$.

\paragraph{Step 7: endpoints.}
Let $z_-$ be an endpoint of $I$, say the left one. The integrated
identity of Step~6 holds for $z_-<z_1<z_2$ in $I^\circ$. As
$z_1\downarrow z_-$, the left side converges because $z\mapsto\E_z\Phi_z$
is continuous on $I$, and the integral converges because $\Psi$ is
continuous on $I$. This also holds if a finite variance vanishes at $z_-$.
In that case,
the corresponding coordinate converges to zero, while the derivative
measures and $\dot K_z$ have continuous limits, and we never invert the
covariance matrix. We obtain
$\E_{z_2}\Phi_{z_2}-\E_{z_-}\Phi_{z_-}=\int_{z_-}^{z_2}\Psi(z)\dd z$,
and the continuity of $\Psi$ gives the one-sided derivative at
$z_-$. We treat the right endpoint and \eqref{eq:fnd-path-normalizer} in the
same way.
\end{proof}
